%% =============================================================================
%%  es.tex — textos de la interfaz en español
%%  Lo carga estilo/estudobiblico.sty cuando \ebIdioma es "es".
%%  Para añadir un idioma: copie este archivo, traduzca los textos y guárdelo
%%  como <codigo>.tex (p. ej. fr.tex); luego use \newcommand\ebIdioma{fr}.
%%  NO cambie las claves (el primer argumento de \EBdeftxt).
%% =============================================================================
% Idioma principal de babel (silabeo, fechas, nombres automáticos)
\def\EB@babelmain{spanish}
% Columna de dados/biblia/livros.tsv con los nombres de los libros
% ("" = columna "nome"; "es" = columna "nome_es"; etc.)
\def\EB@colunalivros{es}

%% ---------------------------------------------------------------- estructura
\EBdeftxt{numeroAbrev}{n.\textsuperscript{o}}
\EBdeftxt{etapa}{Etapa}
\EBdeftxt{parte}{Parte}
\EBdeftxt{sumario}{Índice}
\EBdeftxt{bibliografia}{Bibliografía}
\EBdeftxt{indiceRemissivo}{Índice analítico}
\EBdeftxt{indiceTermos}{Índice de términos originales}
\EBdeftxt{indiceTextos}{Índice de textos bíblicos}

%% --------------------------------------------------------------------- cajas
\EBdeftxt{colofaoUso}{Material de estudio personal.}
\EBdeftxt{colofaoFontes}{Texto griego: SBL Greek New Testament (CC BY 4.0, Society of Biblical Literature \& Logos); morfología MorphGNT (CC BY-SA 3.0). Hebreo y arameo: Westminster Leningrad Codex con morfología de la OpenScriptures Hebrew Bible (CC BY 4.0). Los textos de las traducciones modernas pertenecen a sus titulares y se citan para estudio.}
\EBdeftxt{colofaoTipografia}{Tipografía: EB Garamond, Gentium Book Plus, Marcellus, Source Sans 3, Noto Serif Hebrew y Noto Sans Syriac (SIL Open Font License). Compuesto en \LaTeX\ (LuaLaTeX).}

\EBdeftxt{objetivo}{Objetivo}
\EBdeftxt{nota}{Nota}
\EBdeftxt{pergunta}{Pregunta}
\EBdeftxt{cuidado}{Atención}
\EBdeftxt{aplicacao}{Aplicación}
\EBdeftxt{criticaTextual}{Crítica textual}
\EBdeftxt{baseAT}{Base AT:}
\EBdeftxt{baseNT}{Base NT:}
\EBdeftxt{traducoes}{Traducciones}
\EBdeftxt{noNT}{en el NT}

%% -------------------------------------------------------------- interlinear
\EBdeftxt{grego}{Griego}
\EBdeftxt{hebraico}{Hebreo}
\EBdeftxt{aramaico}{Arameo}
\EBdeftxt{siriaco}{Siríaco}
\EBdeftxt{corpusNT}{En NT}
\EBdeftxt{corpusAT}{En AT}
\EBdeftxt{camadas}{Capas}
\EBdeftxt{original}{original}
\EBdeftxt{transliteracaoMin}{transliteración}
\EBdeftxt{glosaPropria}{glosa propia}
\EBdeftxt{glosaLexical}{glosa léxica (en)}
\EBdeftxt{lemaMin}{lema}
\EBdeftxt{morfologiaMin}{morfología}
\EBdeftxt{varianteMin}{variante}
\EBdeftxt{implicitoMin}{implícito}
\EBdeftxt{variantes}{Variantes}
\EBdeftxt{cabRef}{Ref.}
\EBdeftxt{cabPalavra}{Palabra}
\EBdeftxt{cabTransliteracao}{Transliteración}
\EBdeftxt{cabLema}{Lema}
\EBdeftxt{cabMorfologia}{Morfología}
\EBdeftxt{cabGlosa}{Glosa}
\EBdeftxt{cabStrong}{Strong}
\EBdeftxt{cabNoTrecho}{En el pasaje}
\EBdeftxt{cabGlosasNoTrecho}{Glosas en el pasaje}

%% ----------------------------------------------------------- cuaderno de libro
\EBdeftxt{licao}{Lección}
\EBdeftxt{licoes}{Lecciones}
\EBdeftxt{seloLeia}{LEA}
\EBdeftxt{citado}{citado}
\EBdeftxt{textosBiblicos}{Textos bíblicos}
\EBdeftxt{dicaTextos}{\S\ = punto de la lección · marque al leer el contexto (capítulo) y al comparar traducciones}
\EBdeftxt{semTextos}{(no se encontró ningún texto bíblico en esta lección)}
\EBdeftxt{nTextos}{textos bíblicos}
\EBdeftxt{paraLer}{para leer}
\EBdeftxt{estudadaEm}{Estudiada el}
\EBdeftxt{revisadaEm}{Repasada el}
\EBdeftxt{cabNum}{N\textsuperscript{o}}
\EBdeftxt{cabPag}{Pág.}
\EBdeftxt{cabTextos}{Textos}
\EBdeftxt{cabEstudada}{Estudiada}
\EBdeftxt{cabRevisada}{Repasada}
\EBdeftxt{cabTipo}{Tipo}
\EBdeftxt{cabContexto}{Contexto}
\EBdeftxt{cabTraducoes}{Traducciones}
\EBdeftxt{cabObservacao}{Observación}
\EBdeftxt{cabTexto}{Texto}
\EBdeftxt{cabVezes}{Veces}
\EBdeftxt{cabEmQue}{En qué lecciones}
\EBdeftxt{antesDeLer}{Antes de leer}
\EBdeftxt{antesDeLerDica}{¿Qué pregunta responde esta lección? (con mis palabras)}
\EBdeftxt{ideiaCentral}{Idea central}
\EBdeftxt{ideiaCentralDica}{la respuesta de la lección, en una o dos frases}
\EBdeftxt{anotacoes}{Apuntes}
\EBdeftxt{anotacoesDica}{puntos principales de cada sección; lo que me llamó la atención}
\EBdeftxt{duvidas}{Para investigar}
\EBdeftxt{revisaoSemConsultar}{Repaso sin consultar}
\EBdeftxt{revisaoSemConsultarDica}{responda de memoria a las preguntas de repaso de la lección}
\EBdeftxt{escriturasHebraicas}{Escrituras Hebreas}
\EBdeftxt{escriturasGregas}{Escrituras Griegas}

%% ------------------------------------------------------- cuaderno temático
\EBdeftxt{pesoCentral}{CENTRAL}
\EBdeftxt{pesoApoio}{apoyo}
\EBdeftxt{pesoDificil}{DIFÍCIL}
\EBdeftxt{cabGrupo}{Grupo}
\EBdeftxt{cabPeso}{Peso}
\EBdeftxt{cabLido}{Leído}
\EBdeftxt{cabNota}{Nota}
\EBdeftxt{cabTranslitAbr}{Translit.}
\EBdeftxt{cabOcorr}{Ocurr.}
\EBdeftxt{cabSentido}{Sentido}
\EBdeftxt{siglaAT}{AT}
\EBdeftxt{siglaNT}{NT}
\EBdeftxt{parteDaBiblia}{Parte de la Biblia}
\EBdeftxt{textosEstudados}{Textos estudiados}
\EBdeftxt{ocorrenciasPalavras}{Ocurrencias de las palabras}
\EBdeftxt{base}{Base:}
\EBdeftxt{temaVazio}{(ningún texto con ese criterio en textos.tsv)}
